% !TeX root = ../seminar.tex
% Concept map of the talk, built step by step (the defense's overview frame, redrawn).
\begin{frame}{Overview}
\centering
\begin{tikzpicture}[
    every node/.style={font=\small},
    cbox/.style={terminator, minimum height=2.6em, text width=2.55cm, align=center},
    lab/.style={draw=none, font=\scriptsize\itshape, text=jddarkgrey}]
  \node[cbox, fill=yellow!25] (bh)  at (-5.6, 2.9) {\textbf{Isolated\\ black holes}};
\uncover<1->{
  \node[cbox, fill=yellow!25] (ml)  at (-1.6, 2.9) {\textbf{Gravitational\\ microlensing}};
  \path[connector] (bh) -- node[lab, above, align=center] {only\\ gravity} (ml);
}
\uncover<1->{
  \node[cbox, fill=blue!15] (ph) at (-3.5, 1.15) {Photometry\\ \footnotesize how bright: timescale, parallax};
  \node[cbox, fill=blue!15] (as) at ( 0.3, 1.15) {Astrometry\\ \footnotesize where: the Einstein radius};
  \path[connector] (ml) -- (ph);
  \path[connector] (ml) -- (as);
}
\uncover<1->{
  \node[cbox, fill=blue!15] (m) at (-1.6, -0.6) {\textbf{Lens mass}};
  \path[connector] (ph) -- (m);
  \path[connector] (as) -- (m);
}
\uncover<2->{
  \node[cbox, fill=blue!15] (ro) at (3.6, 2.9) {\textbf{Roman}\\ \footnotesize near-infrared, in space};
  \node[cbox, fill=blue!15] (ru) at (6.5, 2.9) {\textbf{Rubin}\\ \footnotesize optical, ground based};
}
\uncover<2->{
  \node[cbox, fill=green!20] (sim) at (5.05, 1.15) {Monte Carlo\\ simulator};
  \path[connector] (ro) -- (sim);
  \path[connector] (ru) -- (sim);
}
\uncover<2->{
  \node[cbox, fill=green!20] (fi) at (5.05, -0.6) {Fisher forecasts\\ \scriptsize Rubin, Roman, joint};
  \path[connector] (sim) -- (fi);
  \path[connector, dashed] (m) -- (fi);
}
\uncover<3->{
  \node[cbox, fill=red!20, text width=3.4cm] (res) at (1.7, -2.2) {yields \&\\ mass precision};
  \path[connector] (fi) |- (res);
}
\uncover<3->{
  \node[cbox, fill=red!20, text width=3.4cm] (pop) at (-4.0, -2.2) {Black-hole abundance\\ \& mass function};
  \path[connector] (res) -- node[lab, above] {thesis} (pop);
}
\end{tikzpicture}
\note{The seminar takes the first overview report and proposal as its endpoint. Spend about
      3 minutes on the introduction, 8 on microlensing and the Hubble case, 3 on the surveys,
      5 on the simulator and example, 8 on the forecasts and their limits, and 3 on the proposal
      and summary. Use the four photo slides as short transitions. Allow about 20 seconds for the
      animation. Reach the Roman photo near minute 11 and the results near minute 19.}
\end{frame}
